\subsection{Independent Reproduction and Reproducibility Audit}
\label{sec:independent-reproduction}

We conducted an independent CPU-only reproduction to separate values reported
in the source article from measurements produced by our implementation.  The
Dictionary-based Attention (DA) block follows Eqs.~(4)--(6), and DA-HL is solved
from Eqs.~(14)--(16).  The paper does not specify dictionary initialization,
stopping criteria, treatment of non-positive code inner products, the exact
clustering and representation-based hypergraph algorithms, or whether the 100
visual test samples are embedded together with all remaining unlabelled
samples.  We therefore record all additional decisions.  For DA-HL, each
hyperedge has its own dictionary, code similarities are cosine-normalized and
clipped to $[\epsilon,1]$, and the representation hypergraph uses ridge
self-representation followed by selection of the ten largest coefficients. For
the citation experiment, node-centred closed neighborhoods define hyperedges;
a graph-wide dictionary makes the computation scalable, and an attention floor
of 0.5 prevents vertices with disjoint sparse-code support from becoming
disconnected. These are explicit independent implementation choices rather
than undocumented claims about the authors' code.

For MNIST, each of five deterministic repetitions sampled five labelled images
per class and ten test images per class.  The induced graph therefore contained
150 vertices.  We used the reported $K=50$, $\alpha=2^{-4}$, and
$\lambda=0.3$.  Table~\ref{tab:pilot-reproduction} reports mean accuracy and
sample standard deviation across the five repetitions.  DA-HL improved the
$k$-NN hypergraph from $60.40\%$ to $61.60\%$ and the ridge-representation
hypergraph from $36.60\%$ to $42.20\%$, corresponding to gains of 1.20 and 5.60
percentage points.  In contrast, the clustering result decreased from
$48.40\%$ to $47.00\%$.  Thus, this bounded reproduction supports the claim
that dictionary weighting can help noisy local or self-representation
connections, but it does not reproduce a generator-independent improvement.

On Cora, we used the standard Planetoid split of 140 training, 500 validation,
and 1,000 test vertices.  A two-layer incidence propagation network used a
256-dimensional hidden layer, dropout 0.5, Adam, and early stopping.  Across
five initialization seeds, the baseline achieved $81.92\pm1.20\%$ and the DA
variant achieved $81.22\pm1.20\%$.  Although one DA run reached $83.10\%$, the
mean change was $-0.70$ percentage points and therefore did not reproduce the
reported $+1.30$-point Cora gain.

\begin{table}[t]
\centering
\caption{Independent five-seed CPU pilot. Values are mean test accuracy
$\pm$ sample standard deviation. ``Change'' is the paired difference in
percentage points.}
\label{tab:pilot-reproduction}
\scriptsize
\begin{tabular}{lrrr}
\toprule
Dataset / hypergraph & Baseline (\%) & DA (\%) & Change \\
\midrule
MNIST / $k$-NN & $60.40\pm2.88$ & $61.60\pm1.14$ & $+1.20$ \\
MNIST / clustering & $48.40\pm7.67$ & $47.00\pm7.11$ & $-1.40$ \\
MNIST / ridge representation & $36.60\pm3.21$ & $42.20\pm4.71$ & $+5.60$ \\
Cora / incidence network & $81.92\pm1.20$ & $81.22\pm1.20$ & $-0.70$ \\
\bottomrule
\end{tabular}
\end{table}

The discrepancy is not evidence that the published table is incorrect.  It
shows that the published description is insufficient for exact numerical
replication.  In particular, the original processed arrays, random sample
indices, precise hypergraph generators and centroids, dictionary initialization
and convergence settings, handling of negative similarities, and the original
DHGNN code are required to distinguish implementation sensitivity from a true
methodological effect.  We therefore retain the source-paper values as a
separate immutable reference and do not combine them statistically with our
measurements.
